\section{Domain knowledge as evidence for a causal graph}
\label{exec22:domain-scope}

A collections manager considering a payment holiday needs to distinguish a
plausible story from an intervention that improves customer outcomes. A language
model may suggest that financial strain affects delinquency, that servicing
changes financial strain, and that earlier delinquency affects the servicing
offer. Each statement needs a population, a time index, a variable definition
and an evidential basis. Joining the three statements without those qualifications
can create a cycle whose meaning is unclear. A retrieval system can help locate
the underlying evidence, but a retrieved sentence does not by itself establish
an invariant structural mechanism. Two recent clinical studies make the promise
and the limitations of this assistance concrete.

\subsection{Graph elicitation and the meaning of completeness}
\label{exec22:domain-elicitation}

Aziz and Brookhart study a hypothetical comparison of statins with no treatment
for preventing cardiovascular disease. They run four prompting strategies on
GPT-4o and a zero-shot strategy on GPT-o1, with three separate conversations per
configuration. Two author reviewers assess consistency, acyclicity, source
accuracy, completeness and prompt adherence. This is a study of fifteen outputs
in one familiar clinical scenario, rather than an estimate of performance over
a population of banking questions. The original was published on 30 October
2025. \citep{exec22_domain_aziz}

Its completeness criterion asks whether the graph contains seven groups of
variables drawn from an ASCVD risk score. This is a useful check for conspicuous
omissions, but it does not validate the graph's arrows or establish that its
adjustment set identifies the statin effect. Predictors can have different causal
roles: a pretreatment common cause, a mediator and a selection variable cannot
be treated interchangeably simply because each predicts the outcome. In a bank,
a strong arrears predictor measured after a servicing offer may be an outcome
of that offer. Including it in a list is different from correctly deciding
whether to adjust for it.

Table 1 and the article's accompanying text report fourteen acyclic graphs and
seven graphs containing six of the seven criteria. The seven-of-fifteen count
does not support the abstract's ``more than half'' wording. Race is absent from
all fifteen outputs. That is a reproducible observation about the study's
criterion; it does not establish a universal requirement to insert a particular
attribute into every causal model. A social attribute's role needs an explicit
account of the processes it represents and the target population. Blindly
adding a variable to improve a checklist score could conceal the same scientific
problem as omitting it. \citep{exec22_domain_aziz,exec22_domain_reference}

The study also separates producing citations from producing reliable evidence.
One-shot outputs supply sources but include nonexistent references and incorrect
titles, dates or author details. Chain-of-thought outputs give explanations
without sources in all three runs. The GPT-o1 zero-shot configuration consistently
supplies references and code, yet one code output needs manual repair. Its third
run is described as supplying nine sources and then assessed using denominators
of six; the original outputs are needed before aggregating that run's citation
accuracy. The source assessments described in the paper use title and abstract
review, which is weaker than checking that primary methods and results support
a particular causal arrow. \citep{exec22_domain_aziz}

One instruction-based output contains arrows in both directions between statin
use and cholesterol. The authors correctly identify the need for temporally
indexed nodes if feedback is intended. The banking counterpart is
\(D_t\to A_t\to D_{t+1}\), where \(D_t\) is delinquency information available
before action \(A_t\). Collapsing \(D_t\) and \(D_{t+1}\) into a single
``delinquency'' node loses the information clock. A graph assistant should request
that missing distinction rather than silently reverse an arrow to satisfy a
syntactic check. The study supports a brainstorming and source-checking role;
it supplies no validation of a bank effect estimator.

\subsection{Retrieval helps under a particular reference and query}
\label{exec22:domain-retrieval}

Lin and colleagues combine language models, retrieval and observational discovery
for chronic low back pain. Their February 2026 medRxiv original is a preprint,
not a peer-reviewed result. They use baseline observations from 1,868 participants
and a fourteen-variable expert reference. Crucially, the authors explicitly state
that this reference contains cycles and is not a DAG. The reference was developed
from clinical expertise and a literature synthesis; it is a declared comparison
target, not observed intervention truth. \citep{exec22_domain_lin}

The three knowledge systems share Mistral 7B-Instruct v0.3. Ordinary retrieval
supplies document chunks selected by vector similarity. Knowledge-graph retrieval
extracts entities and relations, merges entities, forms communities and retrieves
entities, attached relations, community summaries and supporting text. The paper
describes 600-token chunks with 100-token overlap, an 8,000-token retrieval budget,
and a pipeline involving constrained extraction, a graph database and community
clustering. These are descriptions of the authors' system; they do not establish
that its dependencies are pinned, licensed for the proposed distribution or
operationally integrated here. \citep{exec22_domain_lin}

Both retrieval systems use the same source articles that the experts used to
construct the comparison graph. This is an informative test of reconstructing
an aligned body of knowledge. It does not measure performance when the literature
is incomplete, contradictory, selected by a different team, or independent of the
reference adjudicators. The authors discuss this alignment explicitly. Their
argument that the model predates publication of the expert graph also does not
by itself authenticate the model's training corpus or exclude related earlier
content. The appropriate extension is a deliberately held-out reference and
document design, not a declaration that contamination has been ruled out.

For the data component, the paper describes hill climbing with local generalized
linear models, BIC scores and 500 bootstrap samples. A thresholded consensus graph
forms the initial model. The system then combines its edges with the knowledge
system's proposals and asks a further direction question when both directions
are suggested. Plausibility, association, temporality and a direct causal query
are separate prompt conditions. Their relative performance matters as much as
the strongest result. \citep{exec22_domain_lin}

\begin{center}
\begin{tabular}{p{0.48\linewidth}rr}
\hline
Reported configuration & F1 & Difference \\
\hline
Observational initial model & 0.396 & --- \\
Initial model plus LLM, plausibility & 0.636 & --- \\
Initial model plus ordinary retrieval, plausibility & 0.714 & --- \\
Initial model plus graph retrieval, plausibility & 0.745 & 0.031 over 0.714 \\
Graph retrieval alone after direction resolution & 0.737 & 0.008 below 0.745 \\
\hline
\end{tabular}
\end{center}

These are displayed scores against the single expert graph. The improvement
from 0.396 to 0.745 is 0.349, but the extra gain over ordinary retrieval is 0.031,
and the gain from adding observational discovery to the strongest graph-retrieval
configuration is 0.008. None is a paired confidence interval or a measured return
on engineering effort. Before direction resolution, graph retrieval's plausibility
F1 is 0.602, below the LLM-only value of 0.627. After resolution the values are
0.737 and 0.636. In the final combined systems, ordinary retrieval exceeds graph
retrieval for association by 0.060 and for temporality by 0.130. A blanket ranking
of retrieval architectures would conceal these reversals. \citep{exec22_domain_lin,exec22_domain_reference}

The strongest combined configuration still reports FDR 0.360, meaning disagreement
among its predicted directed edges against this reference under the paper's
definition. It is not a certified probability that a proposed bank causal claim
is false. Small differences also deserve care: subtracting the printed 0.396
from 0.636 gives 0.240, whereas Table 5 prints a change of 0.241. Rounding of
unreported underlying values can explain that discrepancy. It supplies a reason
to recover the original edge records, not a reason to invent corrected results.

\subsection{Direct edges, indirect effects and feedback need different objects}
\label{exec22:domain-semantics}

The chronic-pain plausibility and temporality questions explicitly allow a
relationship ``through one or more intermediate variables.'' Their answers
therefore need not describe direct adjacency among the measured variables.
For the chain
\begin{equation}
 A\longrightarrow M\longrightarrow Y,
 \qquad E=\{(A,M),(M,Y)\},
 \label{exec22:domain-chain}
\end{equation}
\(A\) is an ancestor of \(Y\) even though \((A,Y)\notin E\). A positive
indirect-effect answer cannot be inserted as a direct edge without changing the
graph and possibly its conditional independences. Conversely, when a variable
is unmeasured, a projected graph may legitimately have different edge semantics.
The interface must state which object it is constructing. The paper's source
text is sufficient to identify the mismatch that needs checking; it is not
sufficient to reconstruct every edge's intended semantics. \citep{exec22_domain_lin,exec22_domain_reference}

Taking a union of proposed edges also has no universal accuracy guarantee.
Starting from the correct two-edge set in equation~\eqref{exec22:domain-chain},
adding the false edge \(Y\to A\) produces two true positives, one false positive
and no false negatives. Thus
\begin{equation}
 F_1=\frac{2\,\mathrm{TP}}{2\,\mathrm{TP}+\mathrm{FP}+\mathrm{FN}}
     =\frac45<1.
 \label{exec22:domain-f1}
\end{equation}
It also creates a directed three-cycle without any reciprocal pair. Resolving
only two-way pairs would not remove it. This exact example does not assert that
the paper's code necessarily produces such a failure. It shows why the observed
nonnegative changes in its result table cannot be promoted to a theorem that
combining knowledge and data always improves the graph. \citep{exec22_domain_reference}

There is a further distinction in the observational comparison. Moralization
drops arrowheads and connects parents of a common child. For
\(A\to M\leftarrow B\), it adds the \(A\)--\(B\) link even though that link
is absent from the original skeleton. Consequently, the paper's moralized-graph
comparison is not simply a comparison of adjacencies after ignoring direction.
Good agreement on that object cannot identify a unique causal ordering. Similarly,
a cyclic expert reference and a DAG learner require an explicit account of the
representation mismatch before their distance is interpreted as scientific error.

A useful banking representation retains the proposition type, time horizon,
population, supporting span and evidence strength for each proposed relation.
It distinguishes a direct cause from an ancestor, an association, an ontological
relationship and a modeler's assumption. A type hierarchy or a semantic link in
a knowledge graph is not automatically a causal edge. Hard constraints can narrow
the graph search when justified by an information clock or design rule. Uncertain
literature claims should remain revisable evidence; treating them as inviolable
constraints can exclude the true graph. When sources conflict, returning the
alternatives preserves information that forced orientation destroys.

\subsection{Identical records can support opposite investment decisions}
\label{exec22:domain-value}

An exact illustration explains why better graph agreement alone is insufficient
for a product decision. Let \(A\) denote a centered, standardized servicing
intensity and \(L\) a centered loss index, where one loss unit is stipulated to
equal USD~100. Consider two authored Gaussian structural models. In the first,
all exogenous variables are independent and
\begin{equation}
 A=U_A,\qquad L=-\tfrac12 A+\varepsilon,
 \qquad \operatorname{Var}(U_A)=1,\quad
 \operatorname{Var}(\varepsilon)=\tfrac34.
 \label{exec22:domain-causal}
\end{equation}
In the second, \(U,\eta_A,\eta_L\) are mutually independent centered Gaussians,
each with variance \(1/2\), and
\begin{equation}
 A=U+\eta_A,\qquad L=-U+\eta_L.
 \label{exec22:domain-latent}
\end{equation}
Both models imply exactly the same centered bivariate Gaussian observed law:
\begin{equation}
 \operatorname{Var}\!\begin{pmatrix}A\\L\end{pmatrix}=
 \begin{pmatrix}1&-1/2\\-1/2&1\end{pmatrix}.
 \label{exec22:domain-covariance}
\end{equation}
In the first model, replacing the action equation
by a one-unit higher intervention lowers expected loss by one half unit. In the
second, replacing the action equation leaves \(U\) and \(L\) unchanged, so the
effect is zero. The observed regression slope is \(-1/2\) in both. No amount
of replication of this observed joint law distinguishes the two intervention
effects. The Gaussian loss index is an illustrative local model, not a claim
that realized bank losses can have an unrestricted negative support.
\citep{exec22_domain_reference}

Suppose the higher action costs USD~20 per eligible customer over the declared
horizon, with no other monetized effects. Its modeled incremental margin is
\begin{equation}
 V=-100\{\mathbb E[L\mid\operatorname{do}(A=a+1)]
             -\mathbb E[L\mid\operatorname{do}(A=a)]\}-20,
 \label{exec22:domain-margin}
\end{equation}
which equals USD~30 in the first model and USD~\(-20\) in the second. Retaining
both models therefore retains the range \([-20,30]\). Under a separately declared
subjective weight \(w\) on the first model, expected margin is \(50w-20\),
positive only when \(w>2/5\). A language-model vote fraction, retrieval score or
graph F1 is not that probability. The practical task is to find evidence that
can distinguish the models: an admissible randomization, a justified additional
measurement, or an independently supported restriction. More fluent explanations
of the same observed correlation cannot do so. \citep{exec22_domain_reference}

\subsection{A public evaluation and a concrete integration decision}
\label{exec22:domain-validation}

The proposal should first test graph assistance on public documents and generated
cases with declared causal laws. The evaluation needs two distinct targets.
One is reconstruction of a specified expert account, for which alternative
adjudicators and documented disagreements are legitimate evidence. The other is
recovery of causal implications under a known data-generating law. Conflating
these targets would reward agreement with an expert even when the expert's graph
is incomplete. The Gaussian pair above provides an inexpensive refusal case:
a system given only the shared observed law should retain both effects unless
the task supplies additional identifying information.

Cases should also include indirect-only relations, long cycles, feedback requiring
time expansion, unmeasured common causes, contradictory papers and evidence that
is silent about direction. Keep the adjudication documents separate from the
retrieval corpus for the primary test; report same-corpus reconstruction as a
different condition. Hold out domains as well as individual documents. Match
source access, context allowance and reviewer effort across data-only,
expert-only, ontology-assisted, LLM-only, ordinary-retrieval, graph-retrieval and
combined systems. This tests whether constructing a knowledge graph is worth its
additional extraction, maintenance and review cost.

Measure source identity and span entailment, variable/time alignment, direct-edge
and ancestral accuracy, equivalence-class retention, abstention, valid
identification and decision consequences separately. Repeat prompts and seeds;
do not treat correlated edge votes from a shared corpus as independent experts.
Record graph-building tokens, query tokens, latency, correction effort and
maintenance alongside numerical scores. An 8,000-token context cap does not
equalize the cost of building the two retrieval systems. The public study should
predeclare thresholds for useful assistance and costs before inspecting the
held-out results. There is no measured financial benefit in the two clinical
papers that can set those thresholds for a bank.

For implementation planning, adopt the evidence-provenance pattern and the
separation of variable definition, relation proposal and graph validation.
Integrate retrieval behind a typed interface only after pinning the underlying
corpus, software, model and rights. Retain the ordinary-retrieval comparator;
the results do not justify porting an entire graph-retrieval stack in advance
of a matched benefit test. Reject unconditional union, compulsory orientation
and silent conversion of an indirect relation into a direct edge. The chronic-pain
appendix's Boolean answer schema and fallback to general knowledge lack an
explicit abstention channel; the proposed interface should supply one.
\citep{exec22_domain_lin}

The same reviewed request can pass through a library call, CLI or MCP tool. It
should carry the target, variable clocks, corpus identities, proposed relation
types and permitted operations, and return evidence-backed alternatives and
identification status. Retrieved prose is evidence to inspect, not an instruction
that authorizes another tool call. A cycle or unsupported edge should produce a
specific unresolved result that downstream estimation preserves. These are
proposed contracts, not demonstrated cross-repository integrations. Public
prototype merit must precede an expensive experiment with private bank data.
